% =====================================================================
%  Preámbulo de la tesina
% =====================================================================
\usepackage[T1]{fontenc}
\usepackage[utf8]{inputenc}
\usepackage[english,spanish.tabla.noshorthands.nodecimaldot.noquoting.lcroman]{babel}
\usepackage{lmodern}
\usepackage[final]{microtype}
\usepackage{csquotes}

% --- Página -----------------------------------------------------------
\usepackage[letterpaper,top=2.5cm,bottom=2.5cm,inner=3cm,outer=2.5cm]{geometry}
\setkomafont{disposition}{\normalcolor\bfseries}
\setkomafont{chapter}{\color{udlap}\bfseries\huge}
\setkomafont{section}{\color{udlaposcuro}\bfseries\Large}
\usepackage{scrlayer-scrpage}
\pagestyle{scrheadings}
\clearpairofpagestyles
\ohead{\headmark}
\ofoot*{\pagemark}
\automark[section]{chapter}

% --- Colores UDLAP ------------------------------------------------------
\usepackage{xcolor}
\definecolor{udlap}{HTML}{F37B20}        % naranja del logo
\definecolor{udlaposcuro}{HTML}{9A4A0E}
\definecolor{udlapclaro}{HTML}{FDEBDD}
\definecolor{gris}{HTML}{6B6B6B}
\definecolor{azulfig}{HTML}{2C6FB7}

% --- Matemáticas y unidades -------------------------------------------
\usepackage{amsmath,amssymb,amsthm,mathtools,bm}
\usepackage{esint}
\usepackage{siunitx}
\sisetup{output-decimal-marker={.},group-separator={\,},per-mode=symbol,
         range-phrase={--},list-final-separator={ y }}
\numberwithin{equation}{chapter}
\newcommand{\dd}{\,\mathrm{d}}
\newcommand{\jj}{\mathrm{j}}
\newcommand{\FT}{\mathcal{F}}
\newcommand{\IFT}{\mathcal{F}^{-1}}
\newcommand{\rect}{\operatorname{rect}}
\newcommand{\sinc}{\operatorname{sinc}}
\newcommand{\Real}{\operatorname{Re}}
\newcommand{\Imag}{\operatorname{Im}}
\newcommand{\Dmin}{D_{\min}}
\newcommand{\flim}{f_{\lim}}
\DeclareMathOperator*{\argmax}{arg\,max}

\theoremstyle{definition}
\newtheorem{definicion}{Definición}[chapter]
\theoremstyle{plain}
\newtheorem{proposicion}{Proposición}[chapter]
\renewcommand{\qedsymbol}{$\blacksquare$}

% --- Figuras, tablas, diagramas ---------------------------------------
\usepackage{graphicx}
\graphicspath{{figuras/}}
\usepackage{booktabs,multirow,array,tabularx}
\usepackage[font=small,labelfont={bf,color=udlaposcuro},format=hang]{caption}
\usepackage{subcaption}
\usepackage{tikz}
\usetikzlibrary{calc,arrows.meta,decorations.pathreplacing,decorations.pathmorphing,
                patterns,positioning,shapes.geometric,3d,fit,backgrounds}
\usepackage{tikz-3dplot}
\usepackage{pgfplots}
\pgfplotsset{compat=1.18}
\usepgfplotslibrary{fillbetween}

% --- Algoritmos y código -----------------------------------------------
\usepackage{algorithm}
\usepackage{algpseudocode}
\floatname{algorithm}{Algoritmo}
\renewcommand{\listalgorithmname}{Lista de algoritmos}
\algrenewcommand\algorithmicrequire{\textbf{Entrada:}}
\algrenewcommand\algorithmicensure{\textbf{Salida:}}
\usepackage{listings}
\lstset{basicstyle=\ttfamily\footnotesize,breaklines=true,frame=single,
        rulecolor=\color{gris},keywordstyle=\color{udlaposcuro}\bfseries,
        commentstyle=\color{gris}\itshape,numbers=left,numberstyle=\tiny\color{gris},
        literate={á}{{\'a}}1 {é}{{\'e}}1 {í}{{\'i}}1 {ó}{{\'o}}1 {ú}{{\'u}}1 {ñ}{{\~n}}1}

% --- Cajas ---------------------------------------------------------------
\usepackage[most]{tcolorbox}
% Recuadro para resultados clave
\newtcolorbox{resultado}[1]{colback=white,colframe=udlaposcuro,fonttitle=\bfseries,
  title={#1},breakable}
% Recuadro para cálculos numéricos detallados
\newtcolorbox{calculo}[1][]{colback=gray!6,colframe=gray!60,fonttitle=\bfseries\small,
  title={Cálculo: #1},fontupper=\small,breakable}

% --- Bibliografía APA 7 -------------------------------------------------
\usepackage[style=apa,backend=biber]{biblatex}
\addbibresource{referencias.bib}

% --- Enlaces (al final) -------------------------------------------------
\usepackage[hidelinks,pdfusetitle]{hyperref}
\usepackage[spanish,noabbrev]{cleveref}
\crefname{equation}{la ecuación}{las ecuaciones}
\Crefname{equation}{La ecuación}{Las ecuaciones}
\crefname{figure}{la figura}{las figuras}
\Crefname{figure}{La figura}{Las figuras}
\crefname{table}{la tabla}{las tablas}
\Crefname{table}{La tabla}{Las tablas}
\crefname{chapter}{capítulo}{capítulos}
\Crefname{chapter}{Capítulo}{Capítulos}
\crefname{section}{sección}{secciones}
\Crefname{section}{Sección}{Secciones}
\crefname{subsection}{sección}{secciones}
\crefname{appendix}{apéndice}{apéndices}
\crefname{algorithm}{el algoritmo}{los algoritmos}
\Crefname{algorithm}{El algoritmo}{Los algoritmos}
\crefname{proposicion}{la proposición}{las proposiciones}
\Crefname{proposicion}{La proposición}{Las proposiciones}
